\section{Discussion}
\textbf{What limits each board.} On the Uno the limit is the \MemSramUno\,B SRAM: the state of $n=\MemWallUno$ no longer fits next to \MemOverheadUnoMin--\MemOverheadUnoMax\,B of static use. On the \espA\ it is the free heap, on the \espB\ the largest contiguous block, which is the quantity that matched on this one unit; the \grover\ series on both ESP boards stopped at the per-point time cap before any wall, so whether memory or time ends it beyond that was not tested. The memory wall is consistent with a rule that needs one low-$n$ read (with a predictor chosen per board after the walls were seen), whereas time needs a model; per-gate-kind costs reduce the held-out error from \MapeOne\,\% to \MapeTwo\,\% but remain fits (Section~\ref{sec:results}).

\textbf{Local versus offload.} Offload has a floor of a few milliseconds (serial link, USB latency, host scheduling and the Aer build and run; the split was not tested); local execution on a fast board is faster below the crossover bracket and impossible beyond the wall. Because the PC simulation time grows with $n$ and varies by a median \SimSpreadMed$\times$ (up to \SimSpreadMax$\times$) between sessions on one PC, the brackets of the fast and the slow simulation speed differ and the crossover depends on the host (interpretation: a different PC would move it; not measured), and the bands in Fig.~\ref{fig:f4} should be read as a property of the pair (board, host PC) rather than of the board alone. Only one PC and one baud rate were measured, so how the brackets move with either is not tested.
